
\chapter{Регуляторика: что требует закон}\label{ch:regulation}
Регуляторика выглядит приложением к~технике ровно до момента, когда продукт пытается принять деньги, открыть кошелёк, маршрутизировать транзакции или гарантировать продавцу срок зачисления.
В~этот момент право становится частью архитектуры и задаёт роли участников, разделение ответственности, состояния перевода и обязательные требования к~безопасности.

Глава остаётся в~правовом поле России. Сравнение региональных схем и~стран СНГ -- в~главе~\ref{ch:regional-schemes}. Точные пороги и~даты сверяются по действующим нормативным актам, правилам платёжных систем и актуальным материалам участников рынка.

\section{161-ФЗ: кто именно переводит деньги}\label{sec:regulation-161fz}

161-ФЗ отвечает прежде всего на вопрос: \emph{кто имеет право менять записи} в~реестрах банков.
Без этого словаря сеть, процессинг, банк и платёжный продукт сливаются в~одну фигуру, а~вместе с~этим расплывается ответственность~\cite{fz-161}.

\subsection*{Оператор по переводу денежных средств}

161-ФЗ проводит здесь главное разграничение. Перевод денежных средств выполняет оператор по переводу денежных средств -- банк или НКО, которые вправе принять распоряжение к~исполнению, списать деньги у~плательщика и обеспечить их зачисление получателю~\cite{fz-161}.
У~этих участников -- эмитента, эквайера, банка получателя, кошелька с~клиентским остатком -- возникают денежные обязательства и появляется право менять денежное состояние.

Пока система передаёт сообщения и~остаётся интерфейсом к~чужой денежной функции, правовой вопрос выносится партнёру. Как только продукт ведёт клиентские остатки, инициирует списание от своего имени или обещает пользователю хранение средств, разговор о~статусе должен идти ещё до запуска.

\subsection*{Платёжная система и~её оператор}

Закон отдельно выделяет платёжную систему и оператора платёжной системы. \highlight{Оператор задаёт правила участия, порядок обмена сообщениями, модель расчётов, требования к~бесперебойности и управлению рисками}~\cite{fz-161}.
Жизненный цикл сообщения, обязательные поля, сроки сверки и~правила работы с~инцидентами -- следствия правил конкретной системы. В~карточном мире это видно на примере Visa, Mastercard и Мир.

PSP, платёжный шлюз или оркестратор упрощают интеграцию, но свод правил поверх себя не переписывают. Если схема требует отмену, фиксированный набор полей или обязательную связку с~клиринговой записью, продукт встраивается в~этот свод правил; собственные удобные соглашения его не подменяют.

\subsection*{Три функции платёжной инфраструктуры}

\highlight{161-ФЗ разделяет услуги платёжной инфраструктуры на операционный центр, платёжный клиринговый центр и расчётный центр}~\cite{fz-161}.
Даже если эти роли совмещены у~одного участника, логически это три разных функции. Маршрутизация авторизации, обработка клиринговых записей и~окончательное движение денег работают в~разных слоях инфраструктуры. Сообщение проходит операционный центр, а~расчёт происходит позже и~через другой механизм.

Система, которая маршрутизирует сообщения, не обязана отвечать за расчёт; она работает на уровне сообщения. Продукт, обещающий окончательное зачисление в~определённый момент, опирается на передачу сообщений и на регламент расчёта -- со всеми его сроками и ограничениями. НСПК хорошо показывает это разделение: она держит маршрутизацию и клиринг внутрироссийского карточного трафика, но банки всё равно остаются эмитентами, эквайерами и участниками расчётов.

161-ФЗ намеренно задаёт роли, оставляя реализацию за рамками. Протокол авторизации меняется, обязательство банка перед клиентом~--- нет.
Закон, фиксирующий технологический стек, требовал бы поправки на каждый новый стандарт. 161-ФЗ закрепляет ответственность, а технологическую реализацию оставляет положениям Банка России и~правилам платёжных систем, которые обновляются быстрее закона.

\section{Положения Банка России: как норма становится состоянием системы}\label{sec:regulation-cbr-provisions}

Положения Банка России задают ежедневную работу системы: какие реквизиты обязаны быть в~распоряжении, как защищается платёжная инфраструктура, какие требования возникают к~банковской ИБ, какие документы нужно выдать клиенту по карточной операции.

К~марту 2026 года инженерные требования по ИБ задают действующие 821-П и 851-П~\cite{cbr-821p,cbr-851p}.

\subsection*{762-П: жизненный цикл распоряжения}

762-П -- основной регламент перевода денежных средств. Он определяет реквизиты распоряжения, порядок приёма и исполнения, условия отзыва, прохождение через банки-посредники~\cite{cbr-762p}.
У~любого перевода есть машина состояний: распоряжение получено, принято к~исполнению, исполнено, отклонено, возвращено или отозвано.

762-П раскладывает её на стадии: приём, удостоверение права распоряжения, авторизация, исполнение, уведомление~\cite{cbr-762p}. На примере карточной покупки за 3\,000~руб.:

\begin{enumerate}
  \item \textbf{T+0, 12:00:00}, приём -- банк-эмитент получает
    авторизационный запрос от сети. Момент фиксируется
    в~журнале; с~этой секунды распоряжение считается поступившим.
  \item \textbf{T+0, 12:00:01}, удостоверение права -- хост
    эмитента проверяет PIN или криптограмму, сверяет статус карты.
  \item \textbf{T+0, 12:00:02}, авторизация -- проверка
    достаточности средств, скоринг антифрода, решение
    (approve/отказ).
  \item \textbf{T+1\ldots T+2}, клиринг и исполнение -- эквайер
    передаёт клиринговую запись, сеть рассчитывает нетто-позиции,
    расчётный центр списывает и зачисляет.
  \item \textbf{T+2}, уведомление -- банк-эмитент формирует
    выписку; продавец получает подтверждение зачисления.
\end{enumerate}

Для большинства переводов 762-П устанавливает исполнение в~пределах операционного дня или следующего рабочего дня.

До принятия распоряжения к~исполнению плательщик вправе его отозвать. После начала исполнения -- только пока расчёт не завершён.
Статус \enquote{принято к~исполнению} -- юридическая точка, после которой логика \enquote{отмены} меняется на логику \enquote{возврата}: другие сроки, другие документы, обязательное согласование сторон.

Платёжный API, бэк-офис и внутренняя учётная модель обязаны хранить полную машину состояний перевода с~временн\'{ы}ми метками каждой стадии. Без места для обязательных реквизитов, времени приёма и~времени исполнения система теряет правовой и операционный смысл перевода одновременно.

\subsection*{821-П: защита информации при переводе}

821-П задаёт действующий базовый периметр защиты информации при осуществлении переводов денежных средств~\cite{cbr-821p}. Платёжная инфраструктура -- это контролируемая среда с~требованиями к~защищённым каналам связи, разграничению доступа, журналированию, управлению инцидентами, оценке соответствия и средствам криптографической защиты.

Процессинговый коммутатор и платёжный шлюз не собираются по схеме \enquote{сначала обычный веб-сервис, потом WAF сверху}. Сегментацию, ключи, администрирование и логи планируют на старте, когда архитектура ещё складывается.

Сегментация сети постфактум -- переработка топологии под рабочей нагрузкой; замена алгоритмов криптозащиты после запуска -- миграция данных и~ротация ключей; разграничение привилегий после выдачи десятков ролей -- пересмотр модели доступа. Каждое такое действие дешевле на бумаге, чем в~запущенной системе.

\subsection*{851-П: ИБ банка в~целом}

851-П распространяет требования к~защите информации в~кредитных организациях шире переводов как таковых~\cite{cbr-851p}. Платёжная система встроена в~широкий банковский периметр: каналы (мобильный банк, интернет-банк, ДБО) и процессы (внутренняя разработка, управление уязвимостями, мониторинг и инцидентная отчётность) влияют на неё не меньше, чем сам карточный коммутатор.

Формула \enquote{Мы соответствуем PCI DSS} оставляет вопрос безопасности открытым. Идеально защищённый карточный периметр всё равно живёт внутри регуляторного поля Банка России, где остаются требования к~локальным средствам защиты, процессам и~отчётности.

\subsection*{266-П: карточная операция как документ}

Эмиссию платёжных карт и~операции с~их использованием регулирует Положение Банка России № 266-П от 24.12.2004~\cite{cbr-266p}; обязанность информировать клиента о~каждой операции с~использованием электронного средства платежа дополнительно установлена статьёй~9 161-ФЗ~\cite{fz-161}.
На момент подготовки главы Банк России готовит акт на смену 266-П (на 2024--2026~годах опубликован проект соответствующего Указания); до его вступления в~силу 266-П сохраняет прямое действие.\footnote{По состоянию на май 2026~года заменяющий акт не вступил в~силу; его реквизиты (номер, дата принятия, дата вступления в~силу) проверяйте по первоисточнику на сайте Банка России.}
У~карточного платежа две стороны: авторизационный жизненный цикл и след из документов и уведомлений: чек или иной документ по операции, электронное подтверждение, запись в~мобильном банке, уведомление клиенту, механизм оспаривания.

266-П задаёт состав документа по операции: наименование документа, номер и~дату, сумму и~валюту, вид операции, реквизиты карты, наименование ТСП или банкомата, место проведения -- точный набор зависит от канала и~вида операции~\cite{cbr-266p}.
Принцип один -- каждая операция однозначно идентифицируется по документу, выданному или направленному клиенту.

Минимальный набор реквизитов на уровне каждой транзакции -- исходный материал для ответа клиенту, аудитору и регулятору. Терминальное ПО и страница оплаты формируют и выдают документ с~правильным составом полей. Мобильный банк и~личный кабинет показывают детали операции в~формате, пригодном для оспаривания.

Нормативку полезно читать по уровням. 762-П задаёт жизненный цикл
распоряжения, 821-П и 851-П задают требования к~защите информации,
а 266-П вместе со статьёй~9 161-ФЗ~\cite{cbr-266p,fz-161}
превращают карточную операцию в~документ и клиентский интерфейс.
\section{Лицензирование: какая модель стоит за продуктом}\label{sec:regulation-licensing}

\subsection*{Интерфейс поверх банка-партнёра}

Продукт, который даёт пользователю или продавцу интерфейс, но не берёт обязательства хранить или переводить деньги, обычно строится в~партнёрстве с~банком, эквайером или PSP. На стороне партнёра остаётся денежный статус, расчёты, регуляторная отчётность и формальное исполнение распоряжения.
Сам продукт всё равно попадает в~платёжный периметр и~в~требования по защите информации через договоры и интеграции, но самостоятельным носителем банковской функции не становится.

Архитектура должна рано дать внятные ответы на простые вопросы: где лежат деньги, чьё имя стоит в~договоре с~клиентом, кто выпускает чек, кто отвечает за спорную операцию, кто меняет баланс. Пока ответы не зафиксированы, формула \enquote{мы только интерфейс} не закрывает ни одного из этих вопросов.

\subsection*{Агрегатор, агент и PayFac как пограничная модель}

Между чистым интерфейсом и полноценным хранением денег лежат пограничные модели: платёжный агент, агрегатор, PayFac, иной продукт, который подключает ТСП и управляет клиентским сценарием, хотя денежная функция и расчётная обязанность остаются у~банка или другого партнёра.
Различение проходит по трём признакам: кто держит деньги, кто выпускает расчётный документ, кто входит в~договорные отношения с~клиентом и продавцом.

Подключение ТСП, идентификация клиентов и~ТСП, чек, возврат, спорная операция и расчётные сроки~--- это распределение ролей между участниками из раздела~\ref{sec:ecosystem-participants}: кто подписывает ТСП, кто несёт риск, кто меняет баланс, кто отвечает за расчёт после операции.
Чем активнее продукт берёт на себя ценообразование, маршрутизацию, поддержку клиентов и~операции ТСП, тем раньше нужно зафиксировать, где заканчивается сервис и начинается регулируемая денежная функция. Слова \enquote{агрегатор}, \enquote{платформа} или \enquote{сервис для приёма платежей} ничего не доказывают -- статус определяет фактический набор функций.

\subsection*{Процессинг и маршрутизация без хранения денег}

Платформа, которая маршрутизирует сообщения и~обрабатывает авторизации, но не становится должником клиента по его деньгам, ближе всего к~роли операционного центра. Стоит ей включиться в~подтверждение и~сведение распоряжений~--- и~на неё ложится логика платёжного клирингового центра. В~обоих случаях это услуги платёжной инфраструктуры; деньги остаются у~партнёра~\cite{fz-161}.

Такая система не превращается автоматически в~банк, но быстро упирается в~требования к~бесперебойности, защите информации, аудиту и управлению инцидентами. Формула \enquote{мы только передаём сообщения} снижает лицензионную нагрузку по сравнению с~хранением денег, но не освобождает от инженерной дисциплины.

\subsection*{Кошелёк или баланс клиента}

Продукт, который начинает хранить клиентские средства, вести баланс, открывать кошелёк или иным образом становиться должником пользователя по деньгам, упирается в~статус кредитной организации с~правом на соответствующие операции, включая перевод электронных денежных средств~\cite{fz-161,fz-395-1}.
С~этого момента меняется почти всё: капитал, отчётность, комплаенс, антифрод, контур ПОД/ФТ, безопасность и отношения с~регулятором.

\subsection*{Собственная платёжная система}

Организация, которая сама задаёт правила взаимодействия участников, обязана зарегистрироваться оператором платёжной системы по 161-ФЗ~\cite{fz-161}. Вместе с~регистрацией приходят обязанности по своду правил, управлению рисками, непрерывности работы, допуску участников и модели расчётов.

Пока платформа помогает передавать сообщения, она -- поставщик технологии. Как только она задаёт обязательные правила жизни этих сообщений для нескольких участников, к~ней применяются обязанности оператора платёжной системы.

\practicenote{%
  Главный архитектурный вопрос простой -- возникает ли у~нас собственное
  обязательство по чужим деньгам. Если нет, схема часто строится через
  партнёра; если да, разговор почти неизбежно приходит к~банку, НКО и
  полному регуляторному периметру.

}
\section{Пять вопросов к~любой платёжной архитектуре}\label{sec:regulation-checklist}

Пять вопросов к~новой платёжной схеме или партнёрской интеграции быстро вскрывают правовой, операционный и продуктовый риск.

\paragraph{Список для проверки}
\begin{itemize}
  \item \textbf{Где возникает денежное обязательство?}
    Кто в~этой модели реально становится должником
    пользователя или продавца?
  \item \textbf{Кто вправе принять распоряжение и поменять баланс?}
    Это сразу отделяет интерфейс, процессинг и маршрутизацию
    от банка, НКО или иного участника с~денежной функцией.
  \item \textbf{По какому своду правил живёт операция?}
    Закон, правила платёжной системы, договор с~партнёром
    и внутренняя модель состояний не должны спорить друг с~другом.
  \item \textbf{Какие статусы, документы и уведомления система
    обязана хранить и~выдавать?} Если здесь нет ответа,
    продукт ломается на споре,
    возврате или аудите.
  \item \textbf{В~какой момент операция становится окончательной
    и кто разбирает сбой или спор?} Это вопрос права,
    движения денег, поддержки клиентов
    и продуктового интерфейса одновременно.
\end{itemize}

Фильтр включается прежде всего, когда продукт только рисует верхнеуровневую архитектуру и откладывает определение юридических ролей. Если ответа на вопрос, кто меняет баланс, кто владеет сводом правил, кто разбирает ошибочный платёж, на этом этапе нет, архитектура сырая и концептуально неполная.
Он же включается, когда система собирается в~интеграциях. Если API, договор с~партнёром, сценарий поддержки и внутренняя модель статусов дают разные ответы на одни и~те же пять вопросов, расхождение всплывёт на возврате, споре, аудите или инциденте.

\section{Мир как пример того, как норма попадает в~инженерный план работ}\label{sec:regulation-mir-mandatory}

Обязательная эмиссия и~приём Мир показывают, как регуляторное требование превращается в~инженерные задачи и продуктовые ограничения.

Закон и нормативные акты требуют поддержки Мир -- и~это сразу пункты инженерного плана: корректное распознавание карт Мир, маршрутизация внутрироссийского трафика через инфраструктуру НСПК, нужные сценарии приёма в~терминалах и электронной коммерции, поддержка в~подключении ТСП и~в~модели поддержки~\cite{cbr-ps,cbr-nps-faq}.

План работ PSP или эквайера по обязательной поддержке Мир раскладывается на три трека.
Терминальный трек -- обновление конфигурации TMS, развёртывание ядра бесконтактного приёма Мир, проверка, что POS распознаёт диапазоны BIN карт Мир и корректно обрабатывает контактный и бесконтактный сценарии.
Хостовый трек -- подключение маршрутизации через ОПКЦ, согласование профиля полей ISO~8583 по спецификации НСПК, прохождение хостовой сертификации.
Трек подключения -- включение Мир в~стандартный договор ТСП, настройка сопоставления MCC, корректное отображение в~расчётных отчётах и~на панели управления.

\begin{figure}[tbp]
  \centering
  \includefiguresvg[width=.85\linewidth]{assets/figures/ch06-regulation/regulation-to-engineering}
  \caption{От нормы к~инженерному плану: три трека подключения.}\label{fig:regulation-to-engineering}
\end{figure}

После марта 2022 года поддержка Мир для банков и~процессингов
перестала быть пунктом дорожной карты и~стала отдельной системой --
с~собственной маршрутизацией, терминалами, подключением, поддержкой,
расчётами и~сверкой. Хронология приостановки Visa и Mastercard
и~сложившаяся внутрироссийская конфигурация карточного рынка разобраны в~\ref{sec:mir-history}.

Подробный разбор НСПК, ОПКЦ, MPA, Mir Pay и обязательности Мир -- в~главе~\ref{ch:mir}. Точные пороги обязательного приёма, даты развёртывания и актуальные исключения сверяются по действующим документам НСПК, Банка России и~правилам схем.
